% Extracción quirúrgica del propietario V2 05c: líneas 99--335.
% El resto permanece en este archivo como procedencia, pero no se publica.
\iffalse
\chapter[Macroestructura gauge interna]{Macroestructura gauge interna generada por el continuo discreto}
\label{chap:gauint-macroestructura}

\begin{thesisbox}
La macroestructura de este capítulo se construye únicamente a partir de una
restricción superviviente de \(\mathcal C_{\rm cont}^{\rm disc}\). Los números
\(4\), \(6\) y \(8\), la pantalla, las cuatro involuciones, la medida
proyectiva, el operador de transferencia, el modo de incidencia y el terminal
celular se derivan dentro de esa restricción. Ninguna evaluación posterior
forma parte del dominio ni de las flechas que siguen.
\end{thesisbox}

\section{La semilla proyectiva y la derivación interna de \texorpdfstring{\(4/6/8\)}{4/6/8}}

Sea \(V_{\rm TPK}=\mathbf F_3^2\), con la base ordenada heredada de las dos
hojas del estado TPK. Definimos el conjunto de direcciones
\begin{equation}\label{eq:gauint-cuatro-direcciones}
 I_{\rm gau}:=\mathbf P^1(\mathbf F_3)
 =\{L_\infty,L_0,L_1,L_{-1}\},
\end{equation}
donde
\[
 L_\infty=[1:0],\qquad L_0=[0:1],\qquad
 L_1=[1:1],\qquad L_{-1}=[1:-1].
\]
Por tanto,
\begin{equation}\label{eq:gauint-cardinal-cuatro}
 |I_{\rm gau}|=\frac{3^2-1}{3-1}=4.
\end{equation}
No se introduce un conjunto de cuatro índices: aparece como el cociente de
los ocho vectores no nulos por la acción de \(\mathbf F_3^\times\).

Las incidencias no orientadas son
\begin{equation}\label{eq:gauint-seis-incidencias}
 E_{\rm gau}:=\binom{I_{\rm gau}}2,
 \qquad |E_{\rm gau}|=\binom42=6,
 \qquad Q_{\rm gau}:=\mathbf F_3^{E_{\rm gau}}.
\end{equation}
La pantalla orientada duplica cada dirección, no cada incidencia:
\begin{equation}\label{eq:gauint-ocho-caras}
 \mathscr B_{\rm gau}:=I_{\rm gau}\times\{+,-\},
 \qquad |\mathscr B_{\rm gau}|=2|I_{\rm gau}|=8.
\end{equation}
Así, \(4\), \(6\) y \(8\) proceden respectivamente de direcciones, pares de
direcciones y caras orientadas de una misma semilla.

Para cada \(L\in I_{\rm gau}\), sea \(\vartheta_L\) la involución de la
pantalla que intercambia \((L,+)\) y \((L,-)\) y deja fijas las otras seis
caras. Entonces
\begin{equation}\label{eq:gauint-cuatro-involuciones}
 \vartheta_L^2=I,
 \qquad
 \vartheta_L\vartheta_{L'}=\vartheta_{L'}\vartheta_L
 \quad(L\ne L').
\end{equation}
Las cuatro involuciones se han obtenido antes de definir medida u operador.

\section{Incidencia, modo interno y testigo no constante}

La aplicación de incidencia es
\begin{equation}\label{eq:gauint-incidencia-b}
 B:\mathbf F_3^{I_{\rm gau}}\longrightarrow Q_{\rm gau},
 \qquad
 (Ba)_{\{L,L'\}}:=a_L+a_{L'}.
\end{equation}
Si
\[
 s(q):=\sum_{e\in E_{\rm gau}}q_e\in\mathbf F_3,
\]
cada dirección aparece exactamente tres veces en la suma de las seis
incidencias, y por ello
\begin{equation}\label{eq:gauint-incidencia-invariante}
 s(Ba)=3\sum_{L\in I_{\rm gau}}a_L=0,
 \qquad
 s(q+Ba)=s(q).
\end{equation}

Sobre \(Q_{\rm gau}\) se usa exclusivamente el promedio de conteo
normalizado. El modo de incidencia es la función racional
\begin{equation}\label{eq:gauint-wc-definicion}
 W_c(q):=\mathbf 1_{\{s(q)=0\}}-\frac13.
\end{equation}
Como \(s:Q_{\rm gau}\to\mathbf F_3\) es sobreyectiva, cada fibra tiene
\(3^5\) elementos. De aquí se obtienen, por conteo exacto,
\begin{equation}\label{eq:gauint-wc-momentos}
 \mathbb E_{Q}W_c=0,
 \qquad
 \mathbb E_{Q}W_c^2=\frac29,
 \qquad
 \mathbb E_{Q}W_c^3=\frac{2}{27}.
\end{equation}
Además, \eqref{eq:gauint-incidencia-invariante} da
\begin{equation}\label{eq:gauint-wc-invariancia}
 W_c(q+Ba)=W_c(q).
\end{equation}
Por tanto \(W_c\) es no nulo, centrado e invariante por la incidencia
generada internamente.

\fi
\section{Incidencia, curvatura de acarreo y dominio común de naturalidad}

Para \(x\in\mathcal C_{\rm cont}^{\rm disc}\), denotemos por
\(Y_m(x)\) el conjunto finito de prefijos gauge de profundidad \(m\) que
conservan simultáneamente
\begin{equation*}
\begin{gathered}
 \text{fibra, relaciones de extensión, números, orientación, carry, memoria,}\\
 \text{frontera, ruta, multisección, supervivencia, terminal y lector}.
\end{gathered}
\end{equation*}
Existe una aplicación de collar
\begin{equation}\label{eq:gauint-collar}
 q_m:Y_m(x)\longrightarrow Q_{\rm gau}
\end{equation}
que registra las seis incidencias sin borrar la historia que las produce.
Para \(\ell\ge m\), el truncamiento se escribe
\[
 \tau_{\ell,m}:Y_\ell(x)\longrightarrow Y_m(x).
\]

Para \(y,z\in Y_m(x)\), sea
\begin{equation}\label{eq:gauint-todas-extensiones}
 \operatorname{Ext}^{\rm adm}_{m}(y,z)
\end{equation}
el conjunto de \emph{todas} las extensiones TPK admisibles de \(y\) a \(z\).
Una extensión pertenece a este conjunto sólo si transporta la fibra entera,
las relaciones \(\operatorname{Ext}_\Gamma\), el carry nativo y operatorial,
la orientación, la memoria acumulada, la frontera, la ruta y la condición de
supervivencia. Definimos
\begin{equation}\label{eq:gauint-conteo-kernel}
 a_m(y,z):=\#\operatorname{Ext}^{\rm adm}_m(y,z),
 \qquad
 d_m(y):=\sum_{z\in Y_m(x)}a_m(y,z),
\end{equation}
y, cuando \(d_m(y)>0\),
\begin{equation}\label{eq:gauint-kernel-normalizado}
 k_m(y,z):=\frac{a_m(y,z)}{d_m(y)}.
\end{equation}
Así, \(k_m\) es un conteo normalizado de extensiones completas, no una
función introducida después de proyectar el estado.

La restricción \(\mathcal C_{\rm cont}^{\rm gau}\) está formada por las
historias supervivientes \(x\) para las que se cumplen las condiciones
siguientes a toda profundidad.
\begin{enumerate}
 \item Los conjuntos \(Y_m(x)\) son no vacíos, \(d_m(y)>0\), y los
 truncamientos \(\tau_{\ell,m}\) son sobreyectivos.
 \item Existe una multisección no vacía de pesos racionales estrictamente
 positivos \(\mu_m:Y_m(x)\to\mathbf Q_{>0}\) tal que
 \begin{equation}\label{eq:gauint-medida-proyectiva}
  \sum_{y\in Y_m}\mu_m(y)=1,
  \qquad
  \mu_m(y)=\sum_{\tau_{\ell,m}\widetilde y=y}\mu_\ell(\widetilde y).
 \end{equation}
 \item Cada \(\mu_m\) es estacionaria:
 \begin{equation}\label{eq:gauint-estacionariedad}
  \sum_{y\in Y_m}\mu_m(y)k_m(y,z)=\mu_m(z).
 \end{equation}
 \item La marginal de \((q_m)_*\mu_m\) es el conteo uniforme de
 \(Q_{\rm gau}\). Esta condición hace aplicables literalmente los cálculos
 de \eqref{eq:pdfv2-cor-wc-momentos} al lift \(W_c\circ q_m\).
 \item Se cumple la lumpabilidad directa
 \begin{equation}\label{eq:gauint-lump-directa}
  \sum_{\tau_{\ell,m}\widetilde z=z}
   k_\ell(\widetilde y,\widetilde z)
  =k_m(\tau_{\ell,m}\widetilde y,z)
 \end{equation}
 para todo \(\widetilde y\in Y_\ell\) y \(z\in Y_m\).
 \item Para el kernel adjunto
 \begin{equation}\label{eq:gauint-kernel-adjunto}
  k_m^*(z,y):=\frac{\mu_m(y)k_m(y,z)}{\mu_m(z)},
 \end{equation}
 se cumple la lumpabilidad adjunta
 \begin{equation}\label{eq:gauint-lump-adjunta}
  \sum_{\tau_{\ell,m}\widetilde y=y}
   k_\ell^*(\widetilde z,\widetilde y)
  =k_m^*(\tau_{\ell,m}\widetilde z,y).
 \end{equation}
 \item Las categorías de prefijos poseen la raíz TPK como objeto inicial,
 de modo que su nervio tiene la contracción terminal explícita de
 \eqref{eq:gauint-sdr} más abajo.
\end{enumerate}
Estas condiciones definen el dominio tipado; fuera de él no se publica una
medida o un operador natural por simple declaración.

\section{Medida proyectiva, kernel y naturalidad directa y adjunta}

Fijada una rama de la multisección \(\{\mu_m\}_m\), sea
\begin{equation}\label{eq:gauint-hm}
 H_m:=\operatorname{Fun}(Y_m,\mathbf Q),
 \qquad
 \langle f,g\rangle_m:=\sum_{y\in Y_m}\mu_m(y)f(y)g(y).
\end{equation}
El operador generado por todas las extensiones es
\begin{equation}\label{eq:gauint-km}
 (K_mf)(y):=\sum_{z\in Y_m}k_m(y,z)f(z).
\end{equation}
Las ecuaciones de Markov y estacionariedad dan
\begin{equation}\label{eq:gauint-markov-estacionario}
 K_m\mathbf1=\mathbf1,
 \qquad
 K_m^*\mathbf1=\mathbf1,
 \qquad
 \|K_mf\|_m\le\|f\|_m.
\end{equation}

El lift y la esperanza de truncamiento son
\begin{equation}\label{eq:gauint-lift-esperanza}
 J_{\ell,m}f:=f\circ\tau_{\ell,m},
 \qquad
 E_{m,\ell}:=J_{\ell,m}^*.
\end{equation}
La proyectividad de \(\mu\) implica
\begin{equation}\label{eq:gauint-ej-identidad}
 E_{m,\ell}J_{\ell,m}=I_{H_m}.
\end{equation}
Las dos condiciones de lumpabilidad descargan, respectivamente,
\begin{equation}\label{eq:gauint-naturalidad-directa-adjunta}
 K_\ell J_{\ell,m}=J_{\ell,m}K_m,
 \qquad
 K_\ell^*J_{\ell,m}=J_{\ell,m}K_m^*.
\end{equation}
La segunda igualdad no se infiere de la primera sin la condición adjunta:
ambas quedan registradas en el dominio.

Se define el operador interno positivo
\begin{equation}\label{eq:gauint-transferencia-t}
 T_m:=K_m^*K_m.
\end{equation}
Entonces
\begin{equation}\label{eq:gauint-t-positivo-natural}
 \langle f,T_mf\rangle_m=\|K_mf\|_m^2\ge0,
 \qquad
 T_\ell J_{\ell,m}=J_{\ell,m}T_m.
\end{equation}

\section{Ley común de \(8\) caras y \(4\) cuadrados de Fubini}

Sobre \(Y_m^{\mathscr B_{\rm gau}}\) se define la masa racional
\begin{equation}\label{eq:gauint-ley-ocho-caras}
 \Omega_m^{(8)}(\mathbf y)
 :=\sum_{z\in Y_m}\mu_m(z)
 \prod_{(L,\varepsilon)\in\mathscr B_{\rm gau}}
 k_m(z,y_{L,\varepsilon}).
\end{equation}
Es una probabilidad porque cada factor es Markov. Todas las caras proceden
del mismo estado latente \(z\) y del mismo kernel; no son ocho medidas
independientes añadidas.

Para funciones \(f_{L,+},f_{L,-}\in H_m\), Fubini finito da
\begin{equation}\label{eq:gauint-fubini-ocho}
\begin{aligned}
 &\sum_{\mathbf y}\Omega_m^{(8)}(\mathbf y)
  \prod_{L\in I_{\rm gau}}
  f_{L,+}(y_{L,+})f_{L,-}(y_{L,-})\\
 &\qquad=
 \sum_{z\in Y_m}\mu_m(z)
 \prod_{L\in I_{\rm gau}}
 (K_mf_{L,+})(z)(K_mf_{L,-})(z).
\end{aligned}
\end{equation}
Al tomar \(f_{L,+}=f_{L,-}=f_L\), aparecen simultáneamente los cuatro
cuadrados
\begin{equation}\label{eq:gauint-cuatro-cuadrados}
 \sum_{z\in Y_m}\mu_m(z)
 \prod_{L\in I_{\rm gau}}(K_mf_L(z))^2.
\end{equation}
La marginal de cualquier par opuesto satisface
\begin{equation}\label{eq:gauint-marginal-t}
 \sum_{y_+,y_-}\Omega_{m,L}^{(2)}(y_+,y_-)
 f(y_+)g(y_-)
 =\langle K_mf,K_mg\rangle_m
 =\langle f,T_mg\rangle_m.
\end{equation}
La conmutatividad del producto en \eqref{eq:gauint-ley-ocho-caras} prueba
\begin{equation}\label{eq:gauint-involuciones-ley}
 (\vartheta_L)_*\Omega_m^{(8)}=\Omega_m^{(8)}
 \qquad(L\in I_{\rm gau}).
\end{equation}

\section{Operador interno exacto y testigo propio}

La proyección sobre el sector constante y su complemento son
\begin{equation}\label{eq:gauint-pq}
 P_{\Omega,m}f:=\langle\mathbf1,f\rangle_m\mathbf1,
 \qquad
 Q_m:=I-P_{\Omega,m}.
\end{equation}
El operador celular interno se define por
\begin{equation}\label{eq:gauint-d-hmt}
 D_m^{\rm HMT}:=3Q_m.
\end{equation}
El coeficiente \(3\) es el cardinal del cuerpo de residuos usado en
\eqref{eq:pdfv2-cor-proyectores-cuatro-direcciones}; no representa una
escala añadida.
Para el lift
\[
 W_{c,m}:=W_c\circ q_m
\]
se tiene, por \eqref{eq:pdfv2-cor-wc-momentos},
\begin{equation}\label{eq:gauint-w-eigen}
 P_{\Omega,m}W_{c,m}=0,
 \qquad
 D_m^{\rm HMT}W_{c,m}=3W_{c,m},
 \qquad
 \|W_{c,m}\|_m^2=\frac29.
\end{equation}
Así, el sector complementario posee un testigo no nulo construido por la
incidencia de las seis aristas.

\section{Terminal celular y retracto fuerte}

Sea \(C_*(N\mathsf Y_m;A_n)\) el complejo normalizado del nervio de la
categoría de prefijos, y sea \(y_*\) su raíz inicial. La inclusión y la
aumentación son
\[
 \iota_m:A_n[0]\longrightarrow C_*(N\mathsf Y_m;A_n),
 \quad \iota_m(1)=[y_*],
 \qquad
 \epsilon_m:C_*(N\mathsf Y_m;A_n)\longrightarrow A_n[0].
\]
El objeto inicial proporciona la degeneración extra
\begin{equation}\label{eq:gauint-homotopia-celular}
 H_m[y_0\to\cdots\to y_q]
 :=[y_*\to y_0\to\cdots\to y_q],
\end{equation}
con la convención normalizada cuando aparece una degeneración. Se verifica
\begin{equation}\label{eq:gauint-sdr}
 \epsilon_m\iota_m=I,
 \qquad
 \partial H_m+H_m\partial=I-\iota_m\epsilon_m,
 \qquad
 H_m^2=H_m\iota_m=\epsilon_mH_m=0.
\end{equation}
Este terminal celular es distinto de la proyección constante
\(P_{\Omega,m}\): el primero contrae el nervio y la segunda separa los
sectores del espacio de funciones. Ambos se conservan.

\iffalse
\section{Las trece coordenadas de la restricción gauge}

Para cada \((m,n)\), se define
\begin{equation}\label{eq:gauint-d-trece}
 D_{{\rm gau},m,n}(x):=
 (\mathcal F,\operatorname{Ext}_\Gamma,\operatorname{Rel},
 \mathcal N,\operatorname{or},\operatorname{Carr},\operatorname{Mem},
 \partial,\gamma,\operatorname{Msect},\operatorname{Surv},
 \operatorname{Term},\mathcal L)_{{\rm gau},m,n}(x),
\end{equation}
con el contenido siguiente.
\begin{enumerate}
 \item \(\mathcal F=(V_{\rm TPK},I_{\rm gau},E_{\rm gau},
 Q_{\rm gau},\mathscr B_{\rm gau},Y_m,H_m)\).
 \item \(\operatorname{Ext}_\Gamma\) contiene todos los conjuntos
 \(\operatorname{Ext}^{\rm adm}_m(y,z)\), truncamientos, lifts,
 esperanzas y sus composiciones.
 \item \(\operatorname{Rel}\) contiene \(B\), \(sB=0\), las cuatro
 involuciones, Markov, estacionariedad y las dos lumpabilidades.
 \item \(\mathcal N=(3,4,6,8,m,n,|Y_m|,a_m,d_m)\); esos números no
 sustituyen ninguna de las restantes coordenadas.
 \item \(\operatorname{or}\) conserva las caras \(+/-\), las cuatro
 inversiones \(\vartheta_L\) y la reversión de cada ruta.
 \item \(\operatorname{Carr}\) contiene el carry completo de cada
 extensión contada en \(a_m(y,z)\), incluido el carry ya propagado.
 \item \(\operatorname{Mem}\) conserva el prefijo completo, el estado
 común \(z\) de \(\Omega_m^{(8)}\) y la memoria de sus ocho extensiones.
 \item \(\partial=(q_m,\partial_{N\mathsf Y_m},Q_m,H_m)\) conserva
 frontera de collar, frontera celular y complemento.
 \item \(\gamma\) es el símplice orientado del nervio que registra la ruta
 y no sólo sus extremos.
 \item \(\operatorname{Msect}\) es la familia de todas las medidas
 estacionarias proyectivas compatibles y de todas las extensiones que
 intervienen en los conteos; no se selecciona una rama retrospectiva.
 \item \(\operatorname{Surv}\) registra no vacuidad, profundidad
 arbitraria, proyectividad, las dos lumpabilidades y la persistencia de la
 marginal uniforme de collar.
 \item \(\operatorname{Term}=(\iota_m,\epsilon_m,H_m;
 P_{\Omega,m},Q_m)\) conserva separados los dos terminales.
 \item \(\mathcal L=(\Omega_m^{(8)},K_m,T_m,D_m^{\rm HMT},W_{c,m})\)
 es el lector interno posterior a la construcción de la historia.
\end{enumerate}

Para \(\ell\ge m\) y \(n'\ge n\), todas las coordenadas satisfacen
\begin{equation}\label{eq:gauint-naturalidad-trece}
 u^{\rm gau}_{(\ell,n'),(m,n)}D_{{\rm gau},\ell,n'}(x')
 =D_{{\rm gau},m,n}(\tau_{\ell,m}x').
\end{equation}

\section{Ensamblador, publicador y cadena no ablativa}

Sea
\begin{equation}\label{eq:gauint-graph-d}
 \mathcal G_{\rm gau}:=\operatorname{Graph}(D_{\rm gau}),
 \qquad
 \operatorname{Res}_{\rm gau}(x):=(x,D_{\rm gau}x).
\end{equation}
La multisección dependiente \(\chi_{\rm gau}(x)\) contiene todas las ramas
compatibles de collar, medida y extensión. Definimos
\begin{equation}\label{eq:gauint-x-dependiente}
 X_{\chi,{\rm gau}}
 :=\{(\chi_{\rm gau}(x),x,D_{\rm gau}x):
 x\in\mathcal C_{\rm cont}^{\rm gau}\},
\end{equation}
\begin{equation}\label{eq:gauint-prx}
 \operatorname{pr}_{X,{\rm gau}}(x,D_{\rm gau}x)
 :=(\chi_{\rm gau}(x),x,D_{\rm gau}x).
\end{equation}

El ensamblador crudo
\begin{equation}\label{eq:gauint-ensamblador-crudo}
 a_{\rm gau}:X_{\chi,{\rm gau}}\longrightarrow
 \mathscr M_{\rm gau}^{\rm amb}
\end{equation}
produce la familia completa
\[
 \bigl(I,E,\mathscr B,B,W_c;
 Y_m,\mu_m,k_m,J_{\ell,m},E_{m,\ell};
 K_m,T_m,\Omega_m^{(8)},P_{\Omega,m},Q_m,D_m^{\rm HMT};
 \iota_m,\epsilon_m,H_m\bigr)_{m,n}.
\]
El macroobjeto conserva su entrada mediante
\begin{equation}\label{eq:gauint-graph-a}
 \mathfrak M_{\rm gau}:=\operatorname{Graph}(a_{\rm gau}),
 \qquad
 \mathcal A_{\rm gau}(z):=(z,a_{\rm gau}z).
\end{equation}

El publicador crudo
\begin{equation}\label{eq:gauint-publicador-crudo}
 p_{\rm gau}:\mathfrak M_{\rm gau}\longrightarrow
 \mathscr Y_{\rm gau}
\end{equation}
publica conjuntamente las identidades
\[
 \begin{gathered}
 |I|=4,quad |E|=6,quad |\mathscr B|=8,quad sB=0,
 \quad W_c(q+Ba)=W_c(q),\\
 K_\ell J=JK_m,quad K_\ell^*J=JK_m^*,
 \quad T_m=K_m^*K_m,\\
 (\vartheta_L)_*\Omega_m^{(8)}=\Omega_m^{(8)},
 \quad \text{las cuatro factorizaciones cuadráticas de Fubini},\\
 D_m^{\rm HMT}=3Q_m,quad D_m^{\rm HMT}W_{c,m}=3W_{c,m},
 \quad \partial H_m+H_m\partial=I-\iota_m\epsilon_m.
 \end{gathered}
\]
Finalmente,
\begin{equation}\label{eq:gauint-graph-p}
 \mathcal C_{\rm gau}^{\rm HMT}:=\operatorname{Graph}(p_{\rm gau}),
 \qquad
 \mathcal P_{\rm gau}(M):=(M,p_{\rm gau}M).
\end{equation}
La cadena completa es
\begin{equation}\label{eq:gauint-cadena-completa}
 \mathcal C_{\rm cont}^{\rm disc}\supset
 \mathcal C_{\rm cont}^{\rm gau}
 \xrightarrow{\operatorname{Res}_{\rm gau}}\mathcal G_{\rm gau}
 \xrightarrow{\operatorname{pr}_{X,{\rm gau}}}X_{\chi,{\rm gau}}
 \xrightarrow{\mathcal A_{\rm gau}}\mathfrak M_{\rm gau}
 \xrightarrow{\mathcal P_{\rm gau}}\mathcal C_{\rm gau}^{\rm HMT}.
\end{equation}
Cada flecha tiene una inversa sobre su imagen dada por una proyección de
coordenadas. Ninguna salida reemplaza su historia generadora.

\section{Teorema de realización gauge interna}

\begin{theorem}[Realización interna completa y natural]
\label{thm:gauint-realizacion}
Para todo \(x\in\mathcal C_{\rm cont}^{\rm gau}\), la cadena
\eqref{eq:gauint-cadena-completa} construye, a partir de la estructura
discreta del continuo, una macroestructura natural con cuatro direcciones,
seis incidencias, ocho caras, cuatro involuciones, una medida proyectiva, un
kernel de conteo completo, su transferencia positiva, una ley común de ocho
caras, un testigo de incidencia no constante y un terminal celular. Las
trece coordenadas de \eqref{eq:gauint-d-trece} permanecen recuperables en la
salida.
\end{theorem}

\begin{proof}
Las ecuaciones \eqref{eq:gauint-cuatro-direcciones}--
\eqref{eq:gauint-ocho-caras} derivan \(4/6/8\). La suma de incidencias
\eqref{eq:gauint-incidencia-invariante} prueba la invariancia de \(W_c\), y
el conteo uniforme prueba \eqref{eq:gauint-wc-momentos}. Las extensiones
completas producen \(k_m\) por \eqref{eq:gauint-kernel-normalizado}; las
condiciones explícitas del dominio prueban Markov, estacionariedad y las dos
igualdades de naturalidad. Por tanto \(T_m=K_m^*K_m\) es positivo y natural.
Fubini aplicado una sola vez a \eqref{eq:gauint-ley-ocho-caras} da las cuatro
formas cuadráticas sin cambiar de medida. La marginal uniforme centra
\(W_{c,m}\), de donde \eqref{eq:gauint-w-eigen}. El objeto inicial del nervio
da la degeneración extra y la identidad SDR \eqref{eq:gauint-sdr}. Finalmente,
las tres construcciones por graph retienen sus entradas como primeras
coordenadas, luego la composición no olvida ninguna de las trece entradas.
\end{proof}

\section{Procedencia y falsadores}

\paragraph{Procedencia.}
La semilla TPK, la pantalla orientada, las extensiones supervivientes, la
medida común y la contracción terminal tienen estatuto
\texttt{ARQUITECTURA\_AUTORAL\_PREEXISTENTE}. La transferencia por conteo,
el modo de incidencia y el sector exacto \(3Q\) se presentan como
\texttt{RESULTADO\_RECUPERADO}. La derivación explícita
\(\mathbf P^1(\mathbf F_3)\to4/6/8\), las dos condiciones de lumpabilidad y
el empaquetado no ablativo por graphs tienen estatuto
\texttt{FORMALIZACION\_NUEVA}; formalizan la arquitectura anterior sin
reclamar un origen conceptual nuevo.

\begin{ablation}[Falsadores internos de la rama gauge]
\label{abl:gauint-falsadores}
La construcción queda falsada en el nivel indicado si ocurre cualquiera de
los hechos siguientes.
\begin{enumerate}
 \item Se reemplaza \(I_{\rm gau}\), \(E_{\rm gau}\) o
 \(\mathscr B_{\rm gau}\) por una lista sin los mapas que los derivan.
 \item Falla \(sB=0\) o \(W_c\) deja de ser invariante y no constante.
 \item El conteo \(a_m(y,z)\) omite carry, orientación, memoria, frontera o
 ruta: en tal caso se ha contado otro objeto.
 \item Falla proyectividad, estacionariedad o alguna de las dos
 lumpabilidades; entonces no se puede publicar la naturalidad correspondiente.
 \item Las ocho caras no proceden de una única \(\mu_m\) y un único
 \(k_m\), o alguna \(\vartheta_L\) no preserva \(\Omega_m^{(8)}\).
 \item \(P_{\Omega,m}W_{c,m}\ne0\) o
 \(D_m^{\rm HMT}W_{c,m}\ne3W_{c,m}\).
 \item Falla la identidad SDR \(\partial H+H\partial=I-\iota\epsilon\).
 \item Se omite una de las trece coordenadas: \emph{objeto sustituido;
 conclusión no transferible}.
\end{enumerate}
Estos falsadores controlan pertenencia y naturalidad; no degradan una
historia que satisface todas las identidades.
\end{ablation}
\fi
